\section{Unified Embedding Model for Video, Image, and Visual Document}

Built upon the \vlmvec framework, \model extends embedding learning to a broader range of modalities and tasks, which introduces additional challenges due to their structural and semantic characteristics.
This section presents our methodology, beginning with a brief overview of the VLM2Vec framework (Section 3.1) and the foundation model selection (Section 3.2). Sections 3.3–3.4 then describe the key extensions beyond \vlmvec, including our multimodal data mixing strategy (Section 3.3), and batching strategies to balance data sources during training (Section 3.4).

\subsection{Preliminary: \vlmvec Framework}

As one of the earliest works to adapt MLLMs for universal multimodal embedding learning, \vlmvec~\citep{vlm2vec} formulates multimodal embedding learning as instruction-conditioned embedding alignment across diverse tasks and modalities. In \vlmvec, each training example consists of a query--target pair $(q, t^+)$, where both $q$ and $t^+$ may comprise text, images, videos (extended in \model), or combinations thereof, together with a predefined task-specific instruction. Both $q$ and $t^+$ are independently passed through the same MLLM, and the final embeddings $\mathbf{h}_q$ and $\mathbf{h}_{t^+}$ are obtained from the last token’s hidden state. Training is performed using a standard InfoNCE objective with in-batch negatives and optional hard negatives to align query and target representations, defined as:


\begin{equation}
\mathcal{L}
= -\log 
\frac{\phi(\mathbf{h}_{q}, \mathbf{h}_{t^+})}
{\phi(\mathbf{h}_{q}, \mathbf{h}_{t^+})
+ \sum_{t^- \in \mathcal{N}(q)}
\phi(\mathbf{h}_{q}, \mathbf{h}_{t^-})},
\quad
\phi(\mathbf{h}_q, \mathbf{h}_t)
=
\exp\!\left(\frac{1}{\tau}
\cos(\mathbf{h}_q, \mathbf{h}_t)\right)
\label{eq:infonce}
\end{equation}


where $\cos(\cdot,\cdot)$ denotes the cosine similarity function and $\tau$ is the temperature hyperparameter. Here, $t^-$ denotes a negative target sampled from the in-batch candidates or mined hard negatives associated with the query $q$.
GradCache~\citep{gradcache}, a gradient caching technique, is used to increase the effective batch size for contrastive learning by decoupling backpropagation between the contrastive objective and the encoder, thereby removing encoder backward-pass dependencies along the batch dimension.


\subsection{Foundation Model Selection}

While \vlmvec established an effective paradigm for transforming multimodal large language models into image--text embeddings, its design was primarily tailored to static visual inputs. In contrast, \model extends this paradigm to support interleaved sequences of text, images, and videos, as well as long-form inputs such as multi-page visual documents and full-length videos.

Our goal is to learn a unified embedding space that generalizes across diverse visual modalities and tasks. This necessitates a backbone model that can flexibly encode interleaved multimodal sequences while scaling to long-context inputs. Based on these requirements, we adopt Qwen2-VL as the backbone of \model.

Qwen2-VL is particularly well suited to the challenges posed by \model, offering (i) \texttt{Naive Dynamic Resolution} for efficient processing of inputs with varying spatial resolutions, (ii) \texttt{Multimodal Rotary Position Embedding (M-RoPE)} to capture spatial and temporal structure in videos and visual documents, and (iii) a unified architecture that integrates 2D and 3D convolutions for consistent image and video understanding. Together, these capabilities enable \model to effectively handle diverse and complex multimodal data.



\subsection{Training Data and Multimodal Data Mixing}
% (how we model data pairs of different modalities in a unified way)
To train \model effectively across diverse multimodal tasks, we curate a training dataset comprising three main sources: video-language instruction data, visual document retrieval, and image-based vision tasks.

First, we utilize training data from \texttt{LLaVA-Hound}~\citep{llavahound}, which includes synthetic video-caption pairs and video QA examples generated by ChatGPT. Specifically, we use 300k video-caption pairs and 240k video QA pairs. For the caption data, we adopt two formats: using the caption as the query and video as the target in a video retrieval setup, or using the video as the query to retrieve the most relevant textual description from candidate captions.
We denote these three video training datasets as \texttt{Vid2Cap}, \texttt{Cap2Vid}, and \texttt{VideoQA} in the following discussion.
 
Second, for visual document retrieval tasks, we incorporate datasets from \texttt{ViDoRe}~\citep{colpali} and \texttt{VisRAG}~\citep{visrag}, including \texttt{colpali train set} (118k), \texttt{VisRAG synthetic} (239k), and \texttt{VisRAG in-domain} (123k), which provide training examples for image-based document understanding and retrieval. We denote these three video training datasets as \texttt{ColPali}, \texttt{VisRAG$_{\text{ind}}$}, and \texttt{VisRAG$_{\text{syn}}$} in the following discussion.


Finally, we include image-text datasets from MMEB-train~\citep{vlm2vec} to support generalization across a wide range of visual understanding tasks, including question answering, classification, retrieval, and visual grounding. These datasets help improve the robustness of the learned embeddings across multiple tasks.

We perform on-the-fly batch mixing guided by a pre-defined sampling weight table. This table specifies the relative probabilities of sampling from each dataset, enabling controlled exposure to different task types. By dynamically drawing samples during training, we ensure balanced coverage and prevent overfitting to any single modality or domain. We find that jointly mixing datasets from diverse sources and modalities leads to the most effective model training. Detailed ablation studies are presented in Section~\ref{sec:modality_generalization}. 


% \subsection{Contrastive Learning with Hard-Negative Mining}
% The negative set $\mathcal{N}(q)$ in Eq.~\ref{eq:infonce} is specified according to the task type. For retrieval, ranking, and other general tasks, $\mathcal{N}(q)$ includes both in-batch negatives and mined hard negatives, comprising targets from other examples within the same minibatch as well as hard negatives retrieved or mined for the current query. For classification-style tasks, negatives are constructed within each example: each instance provides one positive target together with a set of task-specific alternatives (e.g., other labels for the same input), and $\mathcal{N}(q)$ is restricted to these in-example negatives. This unified formulation supports effective hard-negative supervision for retrieval while mitigating cross-example false negatives in classification-style settings.

% Hard negatives are mined following the B3 framework~\citep{b3}, which identifies challenging non-positive examples using a pretrained teacher model. For each query, the teacher ranks candidate targets by embedding similarity, and highly ranked candidates that are distinct from the positive are selected as hard negatives and included in $\mathcal{N}(q)$. In this work, VLM2Vec (Qwen2-2B)\footnote{\url{https://huggingface.co/TIGER-Lab/VLM2Vec-Qwen2VL-2B}} is used as the teacher model to provide stable and semantically meaningful similarity estimates. This teacher-guided mining process prioritizes semantically close yet incorrect targets, yielding stronger contrastive supervision than random in-batch sampling while controlling computational overhead.

% }






% Please refer to Appendix~{} for details about the instructions.

% \semih{Add a reference to where the instructions are shared}

% \semih{This subsection so far shows how to obtain $q_{\text{inst}}$. How about the target side formatting?}


\subsection{Batching Strategies}
% (how we use optimized data sampling to improve performance)

To support effective multi-task training over heterogeneous data sources, we introduce an \textit{homogeneous sub-batching} strategy to enhance the hardness and stability of contrastive learning. Specifically, each full batch (e.g., size 1024) is divided into smaller sub-batches (e.g., 8 sub-batches of size 128), where each sub-batch is sampled independently from a single data source only. Compared to per-sample independent sampling, grouping similar samples into sub-batches increases intra-sub-batch homogeneity, which raises the difficulty of contrastive discrimination. At the same time, interleaving multiple such sub-batches preserves cross-task diversity within the full batch, avoiding the instability commonly observed in completely homogeneous batches that originate from a single source. This strategy strikes a balance between sample diversity and structural consistency, fostering more stable and robust optimization dynamics.